\scanpage{6}
\mspage{5}
la forme
\[
 \deleted{\unsure{X}}A_J=\bigcup_{i\in J}A_i^\circ={\unsure{\varphi|\phi}}^{-1}(J)
 \tag*{(1)}
\]
où $\deleted{\unsure{A}}J$ est une partie ouverte de $I$ (ou de $\unsure{\Phi|\Psi|\mathscr{L}}\simeq I$) (\deleted{\unsure{bien}} sûr, cette fois dans $\displaystyle\bigcup_{i\in J}$ il n'est pas question de remplacer $A_i^\circ$ par $A_i$ !). Les parties loc. fermées sont de \ensuremath{=} les
\[
 \deleted{\unsure{X}}A_J=\varphi^{-1}(J)=\bigcup_{i\in J}A_i^\circ
\]
où cette fois $J$ est une partie loc fermée de $I$, i.e. telles que
\[
 i\leq j\leq k,\quad i,k\in J\quad\Rightarrow\quad j\in J
 \tag*{$(*)$}
\]
(les $\deleted{\unsure{J|I}}I_{\geq i}$, $i\in J$, sont des ouvert de $I$ recouvrant $J$, et \deleted{\unsure{dans cet}}\inserted{pour un tel} ouvert, \deleted{$J$}\inserted{$J\cap I_{\geq i}$} est fermé \inserted{dans $I_{\geq i}$,} à cause de $(*)$, donc $J$ est bien loc. fermé si on a $(*)$, l'inv. étant encore plus évident \ldots)

\uline{NB} On peut aussi dire, \inserted{si $\unsure{\mathscr{L}|\mathcal{L}}$ fini,} que les $A^\circ$ $(A\in\mathscr{L})$ sont les \deleted{\unsure{at}} fibres de l'algèbre de Boole engendrée par les $A\in\mathscr{L}$ dans $\unsure{\mathcal{P}|\mathscr{P}}(S)$ $(S=|\mathscr{L}|)$. Les ens obtenus à partir des $A\in\mathscr{L}$, par réunions, intersections et \unsure{complémentaires|complémentations} dans $S=\unsure{\lvert\mathscr{L})|\lvert\mathscr{L}\rvert}$, sont exactement ceux saturés par la rel d'équiv. dans $S$ définie par les $\unsure{A^\circ|X^\circ}$ \unsure{i.e.|id.} celle définie par $\varphi$. Donc ils correspondent biuniv. aux parties quelc. de $I$, par la forme
\[
 J\longmapsto A_J=\varphi^{-1}(J)=\bigcup_{i\in\unsure{J|\supset}}A_i^\circ .
\]
On \unsure{appelera|appellera} \unsure{ses|ces|les} ens \unsure{ainsi|aussi} les \uline{parties} $\mathscr{L}$-\uline{modérées} ou $\mathscr{L}$-\uline{constructibles} de $\mathscr{L}$.
